\begin{abstract}
        本專題使用 FPGA 設計一套聲音錄製、儲存、播放與音量顯示系統。系統透過 Audio CODEC WM8731 進行聲音輸入與輸出，取樣頻率設定為 48 kHz，資料格式為 16 bit。錄音時，聲音資料會先暫存在 SRAM 中，經由使用者確認且解鎖 FLASH 後再寫入 FLASH 作為永久儲存。播放時，系統會先將選定 FLASH slot 的資料載入 SRAM，再由 SRAM 輸出至 Audio CODEC。播放或錄音過程中，系統以 LEDR[17:0] 顯示即時音量條，並以 LCD 與 HEX 顯示模式、狀態、slot、輸入來源與秒數。本專題整合音訊處理、記憶體控制、LCD 顯示、LED 顯示與 FPGA 狀態機設計，達成一套具有互動性的數位音訊應用系統。
\end{abstract}
